\chapter{Approaches to MLC: comparison and dependency map}
\label{ch:approaches-map}

This chapter compares the approaches to MLC that are digested in Chapters~\ref{ch:quadratic-foundations}
and~\ref{ch:digests-foundations}--\ref{ch:digests-sierpinski}, and in the earlier chapters on the
Lyubich renormalization, the Dudko--Lyubich cases, and the frontier.  It adds no new theorem.  For each
approach it records the objects used, the main results and their sources, the conclusion about
the Mandelbrot set that the sources themselves draw, and the formal status in the repository.
The statements are summarized from the digests; the chapters cited give the definitions and the
precise statements.

\section{The approaches in one table}
\label{sec:approaches-table}

\begin{table}[ht]
\centering\small
\begin{tabularx}{\textwidth}{@{}>{\raggedright\arraybackslash}p{0.19\textwidth}>{\raggedright\arraybackslash}X>{\raggedright\arraybackslash}p{0.27\textwidth}>{\raggedright\arraybackslash}p{0.19\textwidth}@{}}
\textbf{Approach} & \textbf{Objects and main results (source)} & \textbf{MLC conclusion stated by the source} & \textbf{Formal status}\\
\hline
Puzzles and parapuzzles &
Dynamical and parameter puzzle pieces; local connectivity of $J$ for non-renormalizable maps (Lyubich notes, Thms~11.5--11.6); Yoccoz's theorems (Lyubich 1993 account) &
MLC at finitely renormalizable parameters; MLC at boundaries of hyperbolic components &
cited; compact-closure criterion proved as a conditional endpoint\\
External rays and B\"ottcher coordinates &
Rays and landing (Milnor, Thms~18.1--18.2); $\M$ connected (Douady--Hubbard; notes, Thm~6.10) &
none by itself &
$K$ connected for $c\in\M$ proved; rays not formalized\\
Hyperbolic components and multipliers &
Multiplier theorem (notes, Thm~6.12); cardioid and bulb; rational analogue (Dudko--Luo, Thm~A) &
MLC at boundaries of hyperbolic components (Yoccoz); $\partial H$ locally connected for rational analogues &
periods~1 and~2 proved, including biholomorphic multiplier maps and real parts\\
Quadratic-like renormalization and straightening &
Quadratic-like maps, $Rf$, hybrid equivalence; straightening (notes, Thm~9.2); codimension-one stability (survey) &
renormalization is the common tool of the later approaches &
interface only\\
A priori bounds &
Definite moduli (Lyubich 1993, Thms~I--II); $W^+_3$ bounds (Dudko--Lyubich, Thms~1.1, 10.1); beau bounds (Feigenbaum Main~A; KKL Theorem~A); Kahn I--III bounds (bounded primitive, decoration, molecule classes); amplification (Kahn IV) &
bounds for all classes, with quasiconformal conjugacy, give MLC (principle in the neutral paper, \S1.5) &
naive lower bounds refuted in the frozen model; genuine bounds not formalized\\
Rigidity &
Combinatorial rigidity (Lyubich 1993, Thm~III and corollary); conformal and affine rigidity of neutral cascades (Dudko--Lim--Lyubich); KKL Mains~B, C &
MLC at the rigid classes, as stated in each source &
interface only\\
Satellite and Molecule renormalization &
Satellite parameters; $\M$ and $J_{c_n}$ locally connected at constructed bounded-satellite $c_n$ (Dudko--Lyubich); molecule condition (Kahn III) &
MLC at the constructed $c_n$ (Dudko--Lyubich); MLC at molecule-condition parameters, conditional on Lyubich's rigidity (Kahn III) &
interface only\\
Feigenbaum points &
Infinitely renormalizable parameters of bounded type; Mains~A, B (Dudko--Lyubich) &
MLC at all such parameters (Main~B, from Main~A and Lyubich's rigidity), including the period-doubling point &
interface only\\
Parabolic primitive renormalization &
Prime types, parabolic tails, parabolically bounded families; Theorems~A--D (KKL) &
MLC pointwise on parabolically bounded families (Theorem~C, conditional on unpublished puzzle input) &
interface only\\
Neutral renormalization &
Pseudo-Siegel bounds, Mother Hedgehog, quasidisks (Dudko--Lyubich, Thms~1.1--1.3); $\psi^\bullet$-maps (Dudko--Luo--Lyubich); bubbles (Dudko--Lim) &
inputs to the satellite programme; the survey states full MLC conditional on the open problems &
none\\
Limbs and combinatorial models &
Homeomorphisms between limbs with equal denominators, preserving periods; canonical lifting (Dudko--Schleicher) &
MLC is equivalent to $\pi\colon\M\to\M_{\rm comb}$ being a homeomorphism (source's reformulation) &
interface only\\
Pacman near Siegel parameters &
Pacman renormalization; unique Siegel pacman; local connectivity of $K_f$ for standard pacmen (Dudko--Lyubich--Selinger, Thms~1.1, 3.12) &
none for $\M$; a later work is announced for satellite parameters &
none\\
Decorations &
Finitely many large components of $\M\setminus\M_s$ (Dudko, decoration theorem) &
none &
none\\
Laminations and matings &
Basilica lamination; mating $\M\cup_{S^1}L_B$ (Dudko, Thms~1.1, 1.2, Cor.~1.3) &
$M_2$ is locally connected; $\M$ is not claimed &
none\\
Core entropy &
$h$ and $\widetilde h$; continuity (Dudko--Schleicher, Thm~1.1); Tiozzo statements (Thm~4.1) &
none &
none\\
Holomorphic motions and hybrid classes &
Extension of holomorphic motions (Sło\-dkowski, Thm~1.3); axiom of choice (Thm~1.4); hybrid classes are analytic of codimension $d(f)$ (Shen, Thm~1.1) &
none in the texts &
none\\
Elephant Eyes &
Class $\mathcal Q(\underline q,\bar b)$; a priori bounds, with no single numbered theorem (Kahn--Lyubich) &
MLC at the corresponding parameters (the paper's claim) &
none\\
Rational analogues &
Sierpi\'nski carpet components of disjoint type; boundedness and $\rho$ (Dudko--Luo, Thms~A--D) &
analogy only &
none\\
Categorical and conditional branch &
Over-approximations in $\mathbf{TopCat}$; conditional roots; refuted frozen certificates (Chapter~\ref{ch:conditional-matrix}) &
conditional MLC root, not an unconditional result &
formal, conditional\\
\end{tabularx}
\caption{Approaches to MLC, the MLC conclusions stated by their sources, and their formal status.}
\label{tab:approaches}
\end{table}

\section{How the approaches depend on one another}
\label{sec:approaches-dependencies}

The sources form a network rather than a list.  The arrows below are the dependencies stated in the sources, in the sense
used in Table~\ref{tab:digest-dependencies}.

\begin{itemize}
\item \emph{Straightening and quadratic-like theory} is the input of Lyubich's 1993 rigidity, of the Dudko--Luo rational
  analogue, and of the Douady--Hubbard identification of hybrid classes.
\item \emph{Yoccoz's puzzle results} are used by Lyubich's 1993 paper, by Dudko's decoration theorem, and by the
  survey's reduction of MLC to rigidity of nested small copies.
\item \emph{A priori bounds} are the common hypothesis of the rigidity statements: the neutral bounds of Dudko--Lyubich feed
  the bubble bounds of Dudko--Lim, the rigidity of neutral cascades (Dudko--Lim--Lyubich), and the $\psi^\bullet$-theory of
  Dudko--Luo--Lyubich.
\item \emph{Pacman renormalization} (Dudko--Lyubich--Selinger) supplies the Siegel and satellite building blocks used by
  the satellite programme, and the survey cites it for the hyperbolicity of Pacman renormalization.
\item \emph{Renormalization of bounded type} gives the Feigenbaum result, and together with puzzle arguments the survey's
  MLC along real lines and veins.
\item \emph{Unpublished input}: the bounded-diameter claim for parabolic tails, used in KKL Main~B, depends on the
  manuscript \emph{Elephant puzzles}; the Dudko--Lyubich neutral bounds depend on earlier work of the same authors.
\item \emph{Combinatorial invariants} (laminations, core entropy, biaccessibility) are analysed in the Douady--Hubbard
  parameter model without assuming MLC; they yield continuity and local connectivity of $M_2$, not of $\M$.
\item \emph{Holomorphic motions} (Słodkowski) are the extension tool of the phase--parameter technique (Chapter~\ref{ch:phase-parameter}),
  and hybrid-class analyticity (Shen) is a complex-analytic counterpart of the straightening theory.
\end{itemize}

\section{What is proved where}
\label{sec:approaches-status}

The repository proves, without assumptions beyond the standard axioms of Lean, the following items from the list above:
the escape criterion, the compactness and exhaustion of outer approximations, the connectedness of filled Julia sets for
$c\in\M$, the real slice $\M\cap\R=[-2,\tfrac14]$, the cardioid and bulb with their multiplier maps and real parts, the
exact examples $c=0$ and $c=-2$, the elementary symmetries of $K_c$, $J_c$ and $\M$, and the obstruction in the frozen Gaussian model.  Every deep analytic theorem in
Table~\ref{tab:approaches} is recorded only as a cited statement or as a Lean interface.  No theorem of the
Yoccoz, Lyubich, Dudko--Lyubich, Kahn--Lyubich, Dudko--Luo--Lyubich, Pacman, Feigenbaum, or parabolic-primitive types is
proved formally in the present repository; for them the Lean development contains only conditional interfaces, and the unconditional statement that $\M$ is locally connected is not formal at all.

\section{Open problems and the remaining gaps}
\label{sec:approaches-open}

Three kinds of gap remain, and they are listed separately so that they are not mistaken for one another.

\begin{enumerate}
\item \emph{Open mathematics.}  The survey's problems (Chapter~\ref{ch:digests-modern}, Section~\ref{sec:digest-survey}) are
  open: uniform bounded geometry of hyperbolic components, uniform a priori bounds for all boundary parameters and all scales,
  a single uniformly hyperbolic renormalization operator, hyperbolicity of neutral and Molecule renormalization, and the
  pseudo-Siegel bounds in the remaining satellite cases.  The rational analogue of the first problem is Theorem~\ref{thm:sierp-A}
  for disjoint type only.
\item \emph{Undigested sources.}  Lyubich's online volumes I--II and the Acta papers, Hu--Jiang, Jiang,
  Inou--Shishikura, Gaidashev--Yampolsky, and Sullivan are cited but not digested in this review.  The Kahn papers (I--IV), the Kahn--Lyubich quasi-additivity and unicritical papers, and Lyubich's Asterisque 261 are digested in Chapters~\ref{ch:digests-apriori-limbs} and~\ref{ch:digests-covering-parapuzzle}.  The Annals paper of Lyubich (1994) is cited but not digested.  Several of the statements
  used above depend on them.  Dudko's papers on canonical decompositions and correspondences are indirectly related and are
  listed in the reference inventory.
\item \emph{Formal gaps.}  The analytic machinery needed for the deep theorems (conformal modulus, quasiconformal surgery,
  straightening, external rays and their landing, and the puzzle theory of Yoccoz) is not in the repository.  The formal
  root remains the conditional statement $\mathrm{MLC.mlc\_conjecture}$, which takes the explicit $\mathrm{RootInput}$
  hypothesis; no global instance of that hypothesis is supplied.
\end{enumerate}

\paragraph{Reading route.}
For the local picture of $\M$, read Chapter~\ref{ch:quadratic-foundations} and Chapter~\ref{ch:digests-monographs}.  For renormalization,
read Chapters~\ref{ch:digests-foundations}, \ref{ch:digests-pacman-neutral}, and \ref{ch:digests-sierpinski}.  For combinatorial
approaches, read Chapter~\ref{ch:digests-laminations-entropy}.  For the status of every claim, read Chapter~\ref{ch:conditional-matrix}
and Appendix~\ref{ch:formal-map}.
